
\subsubsection{4.1 H-E1: Benchmark Independence Verification}

Before testing method-specific differences, the independence of the evaluation benchmarks was verified. The base Llama-2-7B model was evaluated on all three benchmarks, and pairwise Pearson correlations were computed at the item level.

\subsubsection{4.2 H-M1: RLHF Reward Model Training}

A LoRA-based reward model was trained on Llama-2-7B using the HH-RLHF preference dataset. The Bradley-Terry loss learns to predict scalar rewards such that preferred responses receive higher rewards than rejected responses. Training proceeded for 1 epoch on 15,000 preference pairs.

\subsubsection{4.3 H-M2: DPO Training}

DPO was applied to the same base model and dataset. The DPO objective directly optimizes the policy by increasing the log-probability ratio between preferred and rejected responses, weighted by a temperature parameter beta.

\subsubsection{4.4 H-M3: Attractor Analysis}

To test whether different training methods produce distinct behavioral attractors, multiple models were trained with different random seeds (2 seeds per method in quick validation mode). Behavioral embeddings were extracted from 100 probes, and clustering analysis was performed.

\subsubsection{4.5 H-M4: Benchmark Profile Comparison}

Models trained with each method were evaluated on all three benchmarks. Effect sizes (Cohen's d) were computed for each benchmark, and profile correlation was computed to assess whether the methods produce similar or distinct performance patterns across benchmarks.

